%%%%%%%%%%%%%%%%% Оформление ГОСТА%%%%%%%%%%%%%%%%%

% Все параметры указаны в ГОСТЕ на 2021, а именно:

% Шрифт для курсовой Times New Roman, размер – 14 пт.
\setdefaultlanguage[spelling=modern]{russian}
    \setotherlanguage{english}

% Шрифты подключаются явно из ./fonts/ — независимо от системных установок,
% работает одинаково локально и в CI.
% Ligatures=TeX включает: --- → —, -- → –, `` → " и т.д.
\defaultfontfeatures{Path=fonts/, Ligatures=TeX}

\setmainfont{FontsFree-Net-times-new-roman.ttf}
\setromanfont{FontsFree-Net-times-new-roman.ttf}
\newfontfamily\cyrillicfont{FontsFree-Net-times-new-roman.ttf}

% Source Code Pro: загружаем Regular + Bold явно из файла.
% Italic-вариант файла — variable font; пропускаем, чтобы не падать
% на старых сборках TeX Live; курсив синтезируется fontspec'ом при необходимости.
\setmonofont{SourceCodePro-Regular.ttf}[
    BoldFont = SourceCodePro-Bold.ttf,
    AutoFakeSlant = 0.2,
]
\newfontfamily{\cyrillicfonttt}{SourceCodePro-Regular.ttf}[
    BoldFont = SourceCodePro-Bold.ttf,
    AutoFakeSlant = 0.2,
]




% шрифт для URL-ссылок
\urlstyle{same} 

% Междустрочный интервал должен быть равен 1.5 сантиметра.
\linespread{1.5} % междустрочный интервал


% Отступы списков по требованию кафедры:
% «Положение метки 1.25 см, положение текста 0» — метка-нумерация
% стоит как абзацный отступ, тело пункта (включая продолжающиеся
% строки) начинается с левого края страницы, как обычный абзац.
% Это эквивалент Word: «Отступ слева 0, отступ первой строки 1.25 см».
% В enumitem прямой синоним этого — опция wide=<dimen>: пункт
% оформляется как обычный параграф с отступом первой строки.
% Между пунктами — без дополнительного интервала; межстрочный 1.5
% наследуется от \linespread{1.5}.
\setlist[enumerate,1]{
    wide=1.25cm,
    label=\arabic*.,
    labelsep=0.5em,
    topsep=0pt, parsep=0pt, partopsep=0pt, itemsep=0pt,
}
\setlist[itemize,1]{
    wide=1.25cm,
    label=\textendash,
    labelsep=0.5em,
    topsep=0pt, parsep=0pt, partopsep=0pt, itemsep=0pt,
}
\setlist[itemize,2]{label=$\circ$, leftmargin=1cm}



% Каждая новая строка должна начинаться с отступа равного 1.25 сантиметра.
\setlength{\parindent}{1.25cm} % отступ для абзаца

% Текст, который является основным содержанием, должен быть выровнен по ширине по умолчанию включен из-за типа документа в main.tex


%%%%%%%%%%%%%%%%%% Дополнения %%%%%%%%%%%%%%%%%%%%%%%%%%%%%%%%%

% Пути до папок с изображениями
\graphicspath{ {./Images/}{./figures/} }

% Внесение titlepage в учёт счётчика страниц
\makeatletter
\renewenvironment{titlepage} {
	\thispagestyle{empty}
}


% Цвет гиперссылок и цитирования
\usepackage{hyperref} 
 \hypersetup{ 
     colorlinks=true, 
     linkcolor=black, 
     filecolor=blue, 
     citecolor = blue,       
     urlcolor=blue, 
     }
    

% Нумерация рисунков
%\counterwithin{figure}{section}
%\counterwithin{figure}

% Нумерация таблиц
%\counterwithin{table}

%\counterwithin{table}{section}

% шрифт для листингов с лигатурами
% \setmonofont{FiraCode-Regular.otf}[
% 	SizeFeatures={Size=10},
% 	Path = Settings/,
% 	Contextuals=Alternate
% ]
% \setmonofont уже задан выше (см. блок шрифтов с Path=fonts/).
% \setmonofont{Source Code Pro}

% Запрет переносов слов через дефис (по образцу кафедры АСУ).
% Текст «растягивается» межсловными пробелами вместо переноса.
\hyphenpenalty=10000
\exhyphenpenalty=10000
\sloppy
\tolerance=9999
\emergencystretch=4em

% Сквозная нумерация рисунков, таблиц, формул, листингов (по требованиям кафедры).
% При желании заменить на двухуровневую: \counterwithin{figure}{section} и т.д.
% Сейчас оставляем сквозную --- она проще для отслеживания ссылок.

% Запрет на висячие строки (orphans) и заголовков, оторванных от текста (widows).
\widowpenalty=10000
\clubpenalty=10000
\displaywidowpenalty=10000

% Между абзацами не должно быть дополнительного интервала
% (явно фиксируем, чтобы Word-style вставки не привнесли его при копировании).
\setlength{\parskip}{0pt}

% \dottedcontents{<section>}[<left>]{<above-code>}
% {<label width>}{<leader width>}
\dottedcontents{section}[4em]{\bfseries}{4em}{1pc}
\dottedcontents{subsection}[4em]{}{2.5em}{1pc}

\renewcommand{\thesection}{ГЛАВА \arabic{section}}
\renewcommand{\thesubsection}{\arabic{section}.\arabic{subsection}}

% Размеры заголовков по требованию кафедры:
%   - Главы (\section): 16 pt, ЖИРНЫЙ ПРОПИСНЫМИ, по центру
%   - Подразделы (\subsection): 14 pt, жирным (обычный регистр)
%   - Безномерные (\section* — Введение, Заключение): 14 pt, жирным, по центру
% Базовый размер документа — 14 pt, поэтому остальной текст идёт 14 pt.
\usepackage{titlesec}
% В разделе \thesection уже заглавные («ГЛАВА 1»). Применяем
% \MakeUppercase к самому названию через 5-й аргумент \titleformat.
\titleformat{\section}
    {\fontsize{16pt}{19pt}\bfseries\selectfont\centering}
    {\thesection.}{0.5em}{\MakeUppercase}
\titleformat{name=\section,numberless}
    {\fontsize{14pt}{17pt}\bfseries\selectfont\centering}
    {}{0pt}{}
\titleformat{\subsection}
    {\fontsize{14pt}{17pt}\bfseries\selectfont}
    {\thesubsection.}{0.5em}{}
\titlespacing*{\section}      {0pt}{18pt}{12pt}
\titlespacing*{\subsection}   {0pt}{12pt}{6pt}

% настройка подсветки кода и окружения для листингов
%\usemintedstyle{colorful} % делает подсветку для кода
%\newenvironment{code}{\captionsetup{type=listing}}{}


% Посмотреть ещё стили можно тут https://www.overleaf.com/learn/latex/Code_Highlighting_with_minted


% \captionsetup[table]{
%     justification=raggedleft,
%     singlelinecheck=off
% }

% Замена разделителя при цитировании на (,)
\renewcommand*{\multicitedelim}{\addcomma\space}

